\documentclass[10pt]{article}

% -------------------------
% PAQUETES
% -------------------------
\usepackage[utf8]{inputenc}
\usepackage[T1]{fontenc}
\usepackage{geometry}
\usepackage{setspace}
\usepackage{graphicx}
\usepackage{amsmath}
\usepackage{booktabs}
\usepackage{caption}
\usepackage{hyperref}
\usepackage{natbib}
\usepackage{array}
\usepackage{multirow}
\usepackage{listings}  % para insertar el script R
\usepackage{color}

% Opciones de listings para R
\lstset{
  language=R,
  basicstyle=\ttfamily\small,
  breaklines=true,
  frame=single,
  numbers=left,
  numberstyle=\tiny,
  keywordstyle=\color{blue},
  commentstyle=\color{gray},
  showstringspaces=false
}

% -------------------------
% CONFIGURACIÓN DE PÁGINA (según pauta)
% -------------------------
\geometry{
  top=1.5cm,
  left=1.5cm,
  right=1.5cm,
  bottom=2cm
}

\setstretch{1.0}
\setlength{\parindent}{0pt}
\setlength{\parskip}{0.5em}

% -------------------------
% CAPTIONS
% -------------------------
\captionsetup{font=small, labelfont=bf}
\renewcommand{\figurename}{Fig.}
\renewcommand{\tablename}{Tabla}

% -------------------------
% IDENTIFICACIÓN
% -------------------------
\title{\textbf{Control 4 — Análisis de Agrupamiento K-Means sobre Dataset Wine}}
\author{Nataly Huaiquinao, Carlos Ulloa\\
\texttt{nhuaiquinao2021@alu.uct.cl, carlos.ulloa2020@alu.uct.cl}}
\date{\vspace{-0.6em}Control 4, INFO1184, Semestre I-2026}

% -------------------------
\begin{document}

\maketitle
\vspace{-1.0cm}

% -------------------------
% RESUMEN
% -------------------------
\begin{abstract}
Este trabajo aplica el algoritmo K-Means al dataset Wine (178 instancias, 13 atributos químicos)
para determinar el número óptimo de agrupamientos en el rango $k \in [2,15]$.
En la Pregunta~1 se analiza exclusivamente la geometría de los centroides: se calculan las
distancias euclidianas mínimas entre ellos en el espacio estandarizado y se proyectan los datos
mediante PCA para su visualización. La evidencia numérica muestra que en $k=3$ se produce
el primer quiebre significativo de la distancia mínima, los centroides exhiben perfiles químicos
claramente diferenciados y la proyección PCA confirma agrupamientos cohesionados y separados.
La Pregunta~2, que emplea el método del codo (SSE vs.\ $k$), se desarrollará en una etapa
posterior.
\end{abstract}

% -------------------------
% INTRODUCCIÓN
% -------------------------
\section{Introducción}

El dataset Wine reúne resultados de análisis químico de 178 vinos italianos provenientes de
tres cultivares distintos, descritos mediante 13 variables continuas (alcohol, acidez málica,
ceniza, alcalinidad, magnesio, fenoles totales, flavanoides, fenoles no flavanoides,
proantocianinas, intensidad de color, matiz, OD280/OD315 y prolina)
\cite{forina1991wine}.
El objetivo del control es explorar la estructura de agrupamiento latente en estos datos,
sin utilizar las etiquetas originales de clase.

La Pregunta~1 aborda la determinación del número óptimo de clústeres $k$ basándose
\textit{únicamente} en el análisis de centroides: sus valores por variable, las distancias
euclidianas entre ellos y las proyecciones bidimensionales obtenidas mediante Análisis de
Componentes Principales (PCA).
La Pregunta~2 —pendiente— aplicará el método del codo (SSE vs.\ $k$).

El informe se organiza en: Desarrollo (Preguntas~1 y~2), Reflexión, Referencias y Anexo.

% -------------------------
% DESARROLLO
% -------------------------
\section{Desarrollo}

\subsection{Pregunta 1: Determinación del k óptimo mediante análisis de centroides}

\subsubsection{Preprocesamiento}

Los 13 atributos del dataset fueron estandarizados (media~0, desviación~1) antes de aplicar
K-Means, siguiendo la recomendación explícita de \citeauthor{forina1991wine}
(\citeyear{forina1991wine}) de normalizar las variables para algoritmos no invariantes a la
escala \cite{scikit2011}.
El algoritmo se ejecutó con \texttt{n\_init=20} y semilla fija (\texttt{random\_state=42})
para garantizar reproducibilidad.

\subsubsection{Distancias mínimas entre centroides}

Para cada valor de $k$ se calculó la distancia euclidiana mínima entre todos los pares de
centroides en el espacio estandarizado.
Una distancia mínima alta indica que los centroides están bien separados entre sí;
una caída abrupta revela que al agregar un nuevo clúster, al menos dos centroides quedan
muy próximos, evidenciando solapamiento o fragmentación artificial.

\begin{table}[h]
\centering
\small
\begin{tabular}{cccc}
\toprule
$k$ & Dist.\ mín.\ (escalado) & Dist.\ media (escalado) & SSE \\
\midrule
2  & 3.8382 & 3.8382 & 1658.76 \\
\textbf{3}  & \textbf{3.5964} & \textbf{4.1995} & \textbf{1277.93} \\
4  & 2.6674 & 3.9188 & 1175.35 \\
5  & 2.5663 & 3.7957 & 1107.01 \\
6  & 2.5259 & 4.0859 & 1046.00 \\
7  & 2.2928 & 4.4714 &  981.60 \\
8  & 2.2787 & 4.3725 &  935.20 \\
9  & 2.1641 & 4.3686 &  889.89 \\
10 & 1.8279 & 4.3617 &  845.90 \\
11 & 1.8601 & 4.6127 &  825.86 \\
12 & 1.7293 & 4.3725 &  791.30 \\
13 & 1.8862 & 4.4492 &  767.56 \\
14 & 1.7694 & 4.5186 &  742.32 \\
15 & 1.8438 & 4.5515 &  720.50 \\
\bottomrule
\end{tabular}
\caption{Distancias euclidianas entre centroides y SSE para $k = 2\ldots15$.
         Las distancias se calculan en el espacio estandarizado.}
\label{tab:distancias}
\end{table}

La Tabla~\ref{tab:distancias} muestra que la distancia mínima presenta su caída más pronunciada
entre $k=3$ (3.5964) y $k=4$ (2.6674), una reducción de $\approx 0.91$ unidades.
Para $k > 4$ la distancia mínima sigue decreciendo, pero con incrementos de $0.01$–$0.33$
unidades, sin un quiebre comparable.
La Fig.~\ref{fig:dist} ilustra este comportamiento.

\begin{figure}[h]
  \centering
  \includegraphics[width=0.80\textwidth]{fig1_distancia_centroides.png}
  \caption{Distancia euclidiana mínima entre centroides en el espacio estandarizado vs.\ $k$.
           La línea roja señala $k=3$, punto del primer quiebre relevante.}
  \label{fig:dist}
\end{figure}

\subsubsection{Centroides en espacio original para k=3}

La Tabla~\ref{tab:centroides3} muestra los valores medios de cada variable para los tres
centroides obtenidos con $k=3$.

\begin{table}[h]
\centering
\small
\setlength{\tabcolsep}{4pt}
\begin{tabular}{lccccccccccccc}
\toprule
 & Alc. & AcM. & Cen. & AlcCen. & Mg & FenT. & Flav. & FNF & ProAnt & IntCol. & Mat. & OD & Prol. \\
\midrule
C1 & 12.25 & 1.90 & 2.23 & 20.06 & 92.74  & 2.25 & 2.05 & 0.358 & 1.62 & 2.97 & 1.06 & 2.80 & 510  \\
C2 & 13.13 & 3.31 & 2.42 & 21.24 & 98.67  & 1.68 & 0.82 & 0.452 & 1.15 & 7.24 & 0.69 & 1.70 & 619  \\
C3 & 13.68 & 2.00 & 2.47 & 17.46 & 107.97 & 2.85 & 3.00 & 0.292 & 1.92 & 5.45 & 1.07 & 3.16 & 1100 \\
\bottomrule
\end{tabular}
\caption{Centroides de $k=3$ en el espacio original (valores medios por variable).
         Abreviaturas: Alc.=Alcohol, AcM.=Ácido Málico, Cen.=Ceniza, AlcCen.=Alcalinidad Ceniza,
         Mg=Magnesio, FenT.=Fenoles Totales, Flav.=Flavanoides, FNF=Fen.\ No Flavanoides,
         ProAnt=Proantocianinas, IntCol.=Intensidad Color, Mat.=Matiz, OD=OD280/OD315, Prol.=Prolina.}
\label{tab:centroides3}
\end{table}

Las distancias euclidianas entre los tres centroides en el espacio escalado son:

\begin{equation}
  d(C_1, C_2) = 3.9731, \quad d(C_1, C_3) = 3.5964, \quad d(C_2, C_3) = 5.0289
  \label{eq:dist3}
\end{equation}

Los tres valores son mayores a 3.5, lo que confirma una separación geométrica sustancial.
Cada centroide posee un perfil químico diferenciado: $C_1$ presenta baja prolina (510),
baja intensidad de color (2.97) y flavanoides intermedios; $C_2$ muestra alta acidez málica
(3.31), alta intensidad de color (7.24) y muy bajos flavanoides (0.82); $C_3$ se distingue
por el mayor contenido de alcohol (13.68), prolina elevada (1100) y flavanoides altos (3.00).

\subsubsection{Visualizaciones}

La Fig.~\ref{fig:pca3} proyecta los datos y centroides de $k=3$ sobre los dos primeros
componentes principales (varianza explicada: 55.4\%).
Los tres grupos aparecen cohesionados y con escaso solapamiento.

\begin{figure}[h]
  \centering
  \includegraphics[width=0.72\textwidth]{fig2_pca_k3.png}
  \caption{Proyección PCA (PC1 vs.\ PC2) de los datos del dataset Wine con $k=3$.
           Las estrellas negras representan los centroides de cada clúster.}
  \label{fig:pca3}
\end{figure}

La Fig.~\ref{fig:comp} compara las particiones para $k = 2, 3, 4, 5$.
Se observa que $k=2$ fusiona grupos con diferente perfil químico; en $k=4$ y $k=5$ comienzan
a generarse particiones dentro de regiones que ya estaban cohesionadas en $k=3$.

\begin{figure}[h]
  \centering
  \includegraphics[width=0.88\textwidth]{fig3_comparacion_k.png}
  \caption{Proyección PCA para $k \in \{2, 3, 4, 5\}$.
           Las estrellas representan los centroides de cada configuración.}
  \label{fig:comp}
\end{figure}

\subsubsection{Conclusión Pregunta 1}

Con base en el análisis de centroides, el número óptimo de agrupamientos es $\mathbf{k=3}$.
Los argumentos son: (i) la distancia mínima entre centroides experimenta su mayor caída al
pasar de $k=3$ a $k=4$ ($3.5964 \to 2.6674$, diferencia $\approx 0.91$); (ii) los tres
centroides poseen perfiles químicos distintos y bien separados (ec.~\ref{eq:dist3});
(iii) la proyección PCA confirma grupos compactos y no solapados.
Además, $k=3$ coincide con el número real de cultivares del dataset, lo que respalda
la validez de la solución encontrada.

% -------------------------
\subsection{Pregunta 2: Determinación del k óptimo mediante el método del codo}

\textit{Esta sección se completará en una etapa posterior.
La estructura base se conserva para facilitar la incorporación de resultados.}

\subsubsection{Tabla SSE para $k = 2\ldots15$}

\begin{table}[h]
\centering
\begin{tabular}{cc}
\toprule
$k$ & SSE (WCSS) \\
\midrule
\multicolumn{2}{c}{\textit{[Completar con resultados]}} \\
\bottomrule
\end{tabular}
\caption{SSE (Within-Cluster Sum of Squares) para $k = 2\ldots15$.}
\label{tab:sse}
\end{table}

\subsubsection{Gráfica del codo (SSE vs.\ k)}

\begin{figure}[h]
  \centering
  % Reemplazar con la gráfica del codo generada
  \fbox{\parbox{0.6\textwidth}{\centering\vspace{2cm}\textit{[Insertar gráfica SSE vs.\ k]}\vspace{2cm}}}
  \caption{Gráfica del codo: SSE vs.\ número de clústeres $k$.}
  \label{fig:codo}
\end{figure}

\subsubsection{Interpretación}

\textit{[Completar con la identificación del quiebre visual y su justificación numérica.]}

% -------------------------
% REFLEXIÓN
% -------------------------
\section{Reflexión}

El análisis de centroides propuesto en la Pregunta~1 resultó ser un criterio geométrico
robusto: la distancia mínima entre centroides detecta de forma clara el punto en que
agregar un clúster adicional ya no mejora significativamente la separación de los grupos
existentes, sino que fragmenta regiones compactas.
La coincidencia entre el $k=3$ sugerido por este criterio y las tres clases reales del
dataset valida el enfoque, aunque cabe destacar que en un escenario real las etiquetas
no estarían disponibles.

Una limitación del criterio empleado es su sensibilidad a la escala: sin estandarización,
variables como la Prolina (rango $[278, 1680]$) dominarían las distancias y distorsionarían
los resultados.
Como mejora futura, se podría complementar el análisis de centroides con métricas internas
como el índice de Calinski-Harabász, que también se basa en la dispersión entre e
intra-clúster sin requerir etiquetas externas.

% -------------------------
% REFERENCIAS
% -------------------------
\section{Referencias}
\bibliographystyle{apalike}
\bibliography{referencias}

% -------------------------
% ANEXO
% -------------------------
\appendix
\section{Anexo}

\subsection{Script R (archivo insertado)}
% Inserta el script R directamente desde el archivo de trabajo
% ruta relativa desde este archivo: ../Problema.R
\lstinputlisting[language=R]{Problema.R}

\subsection{Bitácora de trabajo}
\begin{itemize}
  \item 2026-04-22: Lectura del enunciado, carga y exploración del dataset Wine.
        Estandarización de variables y ejecución de K-Means para $k=2\ldots15$.
        Cálculo de distancias entre centroides e identificación de $k=3$ como óptimo.
        Generación de figuras PCA y redacción del informe.
  \item \textit{[Completar con sesiones adicionales al incorporar Pregunta 2.]}
\end{itemize}

\end{document}